% !TEX root =./ECE343_Practice_main.tex
%
% L21 practice set (Gauss's Law, Differential Form): the divergence from a
% small box, charge density in Cartesian, cylindrical, and spherical
% coordinates, and the two forms checked against each other.

\fullwidth{\textbf{Learning Objective 5.2: } I can apply the integral and
differential forms of Gauss's Law to calculate the electric field due to charge
density in various geometries, including rectangular, cylindrical, and
spherical coordinate systems.
\\[4pt]\normalfont\small Show your work. Write a \textbf{Documentation} line at the
end (\textbf{None} if you did not collaborate).}

\begin{questions}

%%%%%%%%%%%%%%%%%%%%%%%%%%%%%%%%%%%%%%%%%
\question \textbf{The point form in Cartesian coordinates.}
\begin{parts}

	\part A small box has edges $\Delta x$, $\Delta y$, and $\Delta z$. Find the
	net flux of $\mathbf{D}$ out through its two faces normal to $x$, and show
	that dividing by the volume and letting the box shrink gives
	$\partial D_x/\partial x$.
	\begin{solution}[1.6in]
	The outward normal is $+\xhat$ on the face at $x + \Delta x$ and $-\xhat$
	on the face at $x$, each of area $\Delta y\,\Delta z$:
	\eq{
		\Psi_x &= \lb D_x(x+\Delta x) - D_x(x) \rb \Delta y\, \Delta z \\
		\dfrac{\Psi_x}{\Delta x\,\Delta y\,\Delta z} &= \dfrac{D_x(x+\Delta x) - D_x(x)}{\Delta x}
		\;\xrightarrow{\;\Delta x \to 0\;}\; \dfrac{\partial D_x}{\partial x}
	}
	\end{solution}

	\begin{minipage}[t]{0.58\linewidth}
		\part The flux density in a region is
		$\mathbf{D} = \xhat\,2xy + \yhat\,x^{2} + \zhat\,3z^{2}\ \text{nC/m}^2$,
		with $x$, $y$, $z$ in meters. Find $\rho_v$ at $(1, 2, 3)$.
		\begin{solution}[1in]
		\eq{
			\rho_v = \nabla\mycdot\mathbf{D} &= \dfrac{\partial}{\partial x}(2xy) + \dfrac{\partial}{\partial y}(x^{2}) + \dfrac{\partial}{\partial z}(3z^{2}) \\
			&= 2y + 0 + 6z = 4 + 18 = 22\ \text{nC/m}^3
		}
		\end{solution}
	\end{minipage}\hfill%
	\begin{minipage}[t]{0.37\linewidth}
		\raggedleft \vspace{12pt}
		$\rho_v$ = \ansbox[1.3in]{$22\ \text{nC/m}^3$}
	\end{minipage}

	\part $D_y = x^{2}$ is not zero. Explain why it adds nothing to the
	divergence.
	\begin{solution}[0.8in]
	$D_y$ does not change with $y$, so the flux it carries into a small box
	through one $y$ face leaves through the other. Only a component that
	changes along its own direction produces net flux.
	\end{solution}

\end{parts}

\newpage
%%%%%%%%%%%%%%%%%%%%%%%%%%%%%%%%%%%%%%%%%
\question \textbf{Cylindrical coordinates.} In the region $\rho \le 2\ \m$, the flux
density is $\mathbf{D} = \rhohat\,\rho^{3}/4\ \text{nC/m}^2$, with $\rho$ in meters.
\begin{parts}

	\begin{minipage}[t]{0.58\linewidth}
		\part Find $\rho_v$.
		\begin{solution}[1in]
		Only the radial term of the cylindrical divergence survives:
		\eq{ \rho_v = \dfrac{1}{\rho}\dfrac{\partial}{\partial\rho}\lp\rho\cdot\dfrac{\rho^{3}}{4}\rp = \dfrac{1}{\rho}\cdot\rho^{3} = \rho^{2}\ \text{nC/m}^3 }
		\end{solution}
	\end{minipage}\hfill%
	\begin{minipage}[t]{0.37\linewidth}
		\raggedleft \vspace{12pt}
		$\rho_v$ = \ansbox[1.3in]{$\rho^{2}\ \text{nC/m}^3$}
	\end{minipage}

	\begin{minipage}[t]{0.58\linewidth}
		\part Find the charge per meter of length inside $\rho = 2\ \m$ by
		integrating $\rho_v$.
		\begin{solution}[1in]
		\eq{ \dfrac{Q}{L} = \int_0^{2\pi}\!\!\int_0^{2} \rho^{2}\,\rho\, d\rho\, d\phi = 2\pi\cdot\dfrac{2^{4}}{4} = 8\pi = 25.1\ \text{nC/m} }
		\end{solution}
	\end{minipage}\hfill%
	\begin{minipage}[t]{0.37\linewidth}
		\raggedleft \vspace{12pt}
		$Q/L$ = \ansbox[1.3in]{$8\pi = 25.1$ nC/m}
	\end{minipage}

	\part Check your answer with the integral form of Gauss's law on a
	cylinder of radius $2\ \m$ and length $L$.
	\begin{solution}[1in]
	At $\rho = 2$, $D_\rho = 2^{3}/4 = 2\ \text{nC/m}^2$, and only the side of the
	cylinder carries flux:
	\eq{ \dfrac{1}{L}\oint_S \mathbf{D}\mycdot d\mathbf{s} = D_\rho\cdot 2\pi\rho = 2\times 2\pi\times 2 = 8\pi\ \text{nC/m} }
	The two forms agree.
	\end{solution}

\end{parts}

\newpage
%%%%%%%%%%%%%%%%%%%%%%%%%%%%%%%%%%%%%%%%%
\question \textbf{Spherical coordinates.}
\begin{parts}

	\part Show that the flux density of a point charge,
	$\mathbf{D} = \rhat\,Q/(4\pi r^{2})$, has zero divergence for $r > 0$. Then
	explain how Gauss's law still gives a flux of $Q$ through a sphere around
	the charge.
	\begin{solution}[1.4in]
	\eq{ \nabla\mycdot\mathbf{D} = \dfrac{1}{r^{2}}\dfrac{\partial}{\partial r}\lp r^{2}\cdot\dfrac{Q}{4\pi r^{2}}\rp = \dfrac{1}{r^{2}}\dfrac{\partial}{\partial r}\lp\dfrac{Q}{4\pi}\rp = 0 }
	There is no charge anywhere except at the origin, where the formula fails.
	All of the charge sits at that one point, and any sphere around it encloses
	it, so the flux through the sphere is $Q$.
	\end{solution}

	\begin{minipage}[t]{0.58\linewidth}
		\part A charge distribution produces $\mathbf{D} = \rhat\,k/r$ for
		$r > 0$, where $k$ is a constant. Find $\rho_v$.
		\begin{solution}[1in]
		\eq{ \rho_v = \dfrac{1}{r^{2}}\dfrac{\partial}{\partial r}\lp r^{2}\cdot\dfrac{k}{r}\rp = \dfrac{1}{r^{2}}\dfrac{\partial}{\partial r}\lp kr\rp = \dfrac{k}{r^{2}} }
		\end{solution}
	\end{minipage}\hfill%
	\begin{minipage}[t]{0.37\linewidth}
		\raggedleft \vspace{12pt}
		$\rho_v$ = \ansbox[1.3in]{$k/r^{2}$}
	\end{minipage}

	\part Find the charge inside a sphere of radius $R$ two ways: by
	integrating $\rho_v$, and from the flux of $\mathbf{D}$ through the sphere.
	\begin{solution}[1.5in]
	\eq{
		Q &= \int_0^{2\pi}\!\!\int_0^{\pi}\!\!\int_0^{R} \dfrac{k}{r^{2}}\, r^{2}\sin\theta\, dr\, d\theta\, d\phi = k R\cdot 2\cdot 2\pi = 4\pi k R \\
		Q &= \oint_S \mathbf{D}\mycdot d\mathbf{s} = \dfrac{k}{R}\cdot 4\pi R^{2} = 4\pi k R
	}
	The two counts agree, as the divergence theorem requires.
	\end{solution}

\end{parts}

\vspace{1.5em}
\noindent\textbf{Documentation:} \rule[-2pt]{0.6\linewidth}{0.4pt}

\end{questions}
